\begin{figure}[H]
\centering
\resizebox{\textwidth}{!}{%
\begin{tikzpicture}[node distance=1.5cm and 1.8cm]
    \node[cmsactor] (admin) {Administrador\\del sitio};
    \node[cmssystem, right=of admin] (cms) {Panel CMS\\edición protegida};
    \node[cmsstore, right=of cms] (store) {Almacén de\\contenido};
    \node[cmscontainer, right=of store] (sync) {Sincronización\\de contenido};
    \node[cmssystem, right=of sync] (sitio) {Sitio\\informativo};
    \node[cmsactor, right=of sitio] (visitor) {Visitante\\del sitio};

    \node[cmsboundary, fit=(cms)(store), label={[font=\small]above:Entorno administrativo}] {};
    \node[cmsboundary, fit=(sync)(sitio), label={[font=\small]above:Sitio informativo}] {};

    \draw[cmsarrow] (admin) -- node[above, font=\scriptsize, align=center] {modifica\\contenido} (cms);
    \draw[cmsarrow] (cms) -- node[above, font=\scriptsize, align=center] {guarda\\estructura} (store);
    \draw[cmsarrow] (store) -- node[above, font=\scriptsize, align=center] {copia\\editable} (sync);
    \draw[cmsarrow] (sync) -- node[above, font=\scriptsize, align=center] {actualiza\\datos} (sitio);
    \draw[cmsarrow] (sitio) -- node[above, font=\scriptsize, align=center] {presenta\\información} (visitor);
\end{tikzpicture}%
}
\caption[Arquitectura de contenido editable]{Arquitectura de contenido editable del sitio informativo}
\label{fig:cms_contenido_editable}
\end{figure}
